\phantomsection
\addcontentsline{toc}{chapter}{原课程说明与进度}
\markboth{原课程说明与进度}{}
\begin{center}{\color{structurecolor}\bfseries\fontsize{14.4}{20}\selectfont 原课程说明与进度}\end{center}
本节译自官方 \href{https://ocw.mit.edu/courses/18-152-introduction-to-partial-differential-equations-fall-2011/pages/syllabus/}{教学大纲}与\href{https://ocw.mit.edu/courses/18-152-introduction-to-partial-differential-equations-fall-2011/pages/calendar/}{课程日历}, 不把本整理本的章节顺序解释为原课日历.

课程为 MIT 数学系本科课18.152, 2011年秋季, 教师 Jared Speck. 每周两次课, 每次1.5小时. 先修18.100分析, 以及18.06线性代数、18.700线性代数或18.701代数 I 三者之一. 原课围绕热扩散、静电势和波传播建立方法, 再学习 Fourier 变换、Schrödinger 方程与 Lagrange 场论. 大纲提及 Maxwell 方程等视时间安排的专题; 公开讲义没有建立 Maxwell 系统理论.

指定书为 Sandro Salsa, \textit{Partial Differential Equations in Action: From Modelling to Theory}, Springer, 2010, ISBN9788847007512. 这是商业书目, 不是本项目可再分发的公开原件. 几乎每周一次作业, 通常周四发布、下一周四交, 不接受迟交, 最低一次作业成绩剔除. 期中90分钟, 期末综合全课且稍侧重后半段; 大纲称均闭卷、不得用笔记; 实际期末试卷允许一张正反面手写笔记, 见期末章说明, 以原卷为准. 总评作业30\%, 期中30\%, 期末40\%. 原日历仅列课序, 未列逐课公历日期; 下表不推测日期.

\begin{longtable}{p{.12\textwidth}p{.58\textwidth}p{.21\textwidth}}
\toprule 课序&官方主题的中文译名&提交/考试\\\midrule\endhead
L1&偏微分方程导论&\\
L2&热方程导论&\\
L3&热方程唯一性&作业1\\
L4&热方程弱最大原理及基本解导论&\\
L5&热基本解与全空间 Cauchy 问题&作业2\\
L6&Laplace 与 Poisson 方程&\\
L7&Poisson 基本解&作业3\\
L8&Poisson 的 Green 函数&\\
L9&Poisson 公式、Harnack 不等式、Liouville 定理&作业4\\
L10&波方程导论&作业5\\
L11&波方程球面平均法&\\
L12&Kirchhoff 公式与 Minkowski 几何&作业6\\
L13&波方程几何能量估计&\\
E1&期中考试&90分钟\\
L14&波方程几何能量估计续&\\
L15&二阶方程分类&作业7\\
L16&Fourier 变换导论&\\
L17&Fourier 变换续&作业8\\
L18&Fourier 反演与 Plancherel 定理&\\
L19&Schrödinger 方程导论&作业9\\
L20&Schrödinger 方程续&\\
L21&Lagrange 场论导论&选做 Bonus\\
L22&Lagrange 场论续&作业10\\
L23&Lagrange 场论续&\\
L24&输运与 Burgers 方程&作业11\\
E2&期末考试&综合全课\\\bottomrule
\end{longtable}
